%% kapa: exam=1 round=20 year=2009 period=2 form=A subject=appraisal_law subject_ko=부동산관계법규 seq=8 number=48
%% source: Q-Net 감정평가사 기출문제 (cst00309) artlSeq=5024178 file=extracted/1008/5024178/2009년 제20회 감정평가사 1차시험 기출문제/2교시.hwp
%% conversion: kordoc 4.24.51 -> tools/gen1.py (auto); review=proofread-sonnet; flags=none
\begin{question}{48}
국토의 계획 및 이용에 관한 법령상 용도지역 및 용도지구 안에서의 건축제한에 관한 설명으로 옳지 \uline{않은} 것은?
\begin{choices}
  \item 자연취락지구 안에서는 도시계획조례로 10층 이하의 공동주택의 건축을 허용할 수 있다.
  \item 고도지구 안에서는 도시관리계획으로 정하는 높이를 초과하거나 미달하는 건축물을 건축할 수 없다.
  \item 경관지구 안에서의 건축물의 건폐율, 용적률, 높이 등에 관하여는 그 지구의 경관의 보호ㆍ형성에 필요한 범위 안에서 도시계획조례로 정한다.
  \item 특정용도제한지구 안에서는 주거기능을 훼손하거나 청소년 정서에 유해하다고 인정하여 도시계획조례가 정하는 건축물을 건축할 수 없다.
  \item 자연녹지지역에서는 지목이 대(垈)라 하더라도 아파트를 건축할 수 없다.
\end{choices}
\end{question}
\answer{1}
